\section{\textit{Internet of Things} dan Keterbatasan Sumber Daya Perangkat \textit{Edge}}

\textit{Internet of Things} (IoT) merupakan jaringan objek fisik, seperti sensor, aktuator, dan perangkat pintar, yang saling terhubung dan mampu bertukar data melalui internet tanpa interaksi manusia secara langsung.
Jumlah perangkat IoT diproyeksikan melampaui 60 miliar unit pada 2030 dengan karakteristik yang heterogen dari sisi kemampuan komputasi, energi, dan konteks penggunaan \parencite{ahsan_comprehensive_2025}.

Untuk mengelola kompleksitas tersebut, arsitektur aplikasi IoT umumnya diklasifikasikan menjadi \textit{cloud, edge,} dan \textit{fog computing} \parencite{dauda_survey_2024}.
\textit{Cloud} menyediakan pemrosesan dan penyimpanan terpusat untuk skala besar, tetapi dapat menghadapi latensi dan keterbatasan \textit{bandwidth}.
\textit{Edge computing} memindahkan komputasi lebih dekat ke sumber data untuk mendukung pemrosesan \textit{real-time} dan mengurangi beban jaringan, sedangkan \textit{fog computing} mengombinasikan karakteristik \textit{cloud} dan \textit{edge} untuk lingkungan IoT yang lebih kompleks.

Pemindahan komputasi ke \textit{edge} tidak menghilangkan keterbatasan sumber daya, tetapi memindahkannya ke perangkat tempat komputasi dijalankan.
Perangkat IoT pada \textit{edge} umumnya memiliki keterbatasan memori, daya komputasi, dan energi, sehingga mekanisme yang dijalankan perlu dirancang ringan agar tidak bergantung sepenuhnya pada komputasi awan \parencite{canavese_security_2024}.

Kebutuhan akan algoritma ringan dan adaptif telah mulai direspons dalam optimasi metaheuristik untuk \textit{edge computing}.
Salah satu contohnya adalah \textit{Modified Greylag Goose Optimization} (MGGO) yang dirancang untuk alokasi layanan IoT pada \textit{edge} dengan mekanisme adaptif tanpa pelatihan atau pemodelan eksplisit yang berat \parencite{khosrowshahi_refined_2025}.
Hal ini menunjukkan bahwa keterbatasan sumber daya \textit{edge} dapat ditangani melalui desain algoritma yang sesuai, bukan hanya melalui peningkatan kapasitas perangkat keras.

Keterbatasan memori, komputasi, dan energi tersebut menjadi premis dasar penelitian ini dalam merancang algoritma metaheuristik yang ringan untuk tugas \textit{policy mining} ABAC pada perangkat \textit{edge-IoT}.